\begin{corollary}[极小 prime-slice 编码同时刚性化偏序、原子支撑与块内配对计数]\label{cor:conclusion-partition-primeslice-divisibility-valuation-paircount}
在定理 \ref{thm:conclusion-partition-primeslice-minimal-uniqueness} 的设定下，对任意 \(\pi,\sigma\in\mathcal P([n])\) 都有
\[
\pi\le \sigma
\iff
\Phi(\pi)\mid \Phi(\sigma).
\]
并且对任意边 \(e\in E_n\)，有
\[
v_{\eta(e)}\bigl(\Phi(\pi)\bigr)=\mathbf 1_{\{e\in F_\pi\}},
\]
以及
\[
\Omega\bigl(\Phi(\pi)\bigr)
=
|F_\pi|
=
\sum_{B\in\pi}\binom{|B|}{2}.
\]
因此，在极小素数预算下，分割格的偏序、原子支撑与块内配对总数都被同一个平方自由整数 \(\Phi(\pi)\) 完全而唯一地承载。
\end{corollary}
\begin{proof}
由定理 \ref{thm:conclusion-partition-primeslice-minimal-uniqueness}，
\[
\Phi(\pi)=\prod_{e\in F_\pi}\eta(e),
\qquad
\Phi(\sigma)=\prod_{e\in F_\sigma}\eta(e).
\]
由于 \(\eta(e)\) 两两互异且 \(\Phi(\pi),\Phi(\sigma)\) 平方自由，故
\[
\Phi(\pi)\mid \Phi(\sigma)
\iff
F_\pi\subseteq F_\sigma.
\]
而正文中的 partition-closure 判据给出
\[
F_\pi\subseteq F_\sigma
\iff
\pi\le \sigma.
\]
这就得到第一式。

第二式同样由平方自由分解立即得出：素数 \(\eta(e)\) 在 \(\Phi(\pi)\) 中出现当且仅当 \(e\in F_\pi\)，且其指数恰为 \(1\)。最后，因为 \(\Phi(\pi)\) 平方自由，\(\Omega(\Phi(\pi))\) 正是其不同素因子个数，故
\[
\Omega\bigl(\Phi(\pi)\bigr)=|F_\pi|.
\]
而
\[
F_\pi=\bigsqcup_{B\in\pi}\binom{B}{2},
\]
故
\[
|F_\pi|=\sum_{B\in\pi}\binom{|B|}{2}.
\]
证毕。
\end{proof}
